\vspace{-8pt}
\section{Limitations and Discussion}
\vspace{-3pt}

In this paper, we present an orthogonal approach to improve the performance of language models using multi-agent debate. We find that the approach is simple and effective across a wide set of different reasoning and validity language modeling tasks.

\myparagraph{Limitations.} In comparison to other prompting techniques, our multiagent debate procedure is more computationally expensive, as it requires both multiple language generations, and an underlying debate procedure. However, we believe that this approach may be seen as a method to generate additional data that may be distilled back to self-improve the original base model.

Further, we observed that as debates became longer in duration, current language models sometimes struggled to fully process the entire debate input, and typically only focused on the most recent generations. We believe that this performance will be alleviated with longer-context and improved language models or by summarizing early portions of the debate.

Finally, we found that while debates typically converged into single final answers, these answers were not necessarily correct. Despite answers being incorrect, language models would confidently affirm that their answer is correct and consistent with all other agent responses. We believe this result is in part due to the fact that LMs do not correctly express their uncertainty when generating responses, and believe that other orthogonal approaches to improve this performance would improve the results of multiagent debate.